% =====================================================================
% ozGUI / ozLib -- MASTER report.
%
% This file is the single source of truth for the combined report. It pulls
% in the two documents via \input; docmute strips everything outside their
% \begin{document} ... \end{document}, so
%
%   * the two documents remain INDEPENDENTLY BUILDABLE, exactly as before
%     -- their own preambles are untouched and are used when they are
%     compiled on their own;
%   * this master supplies the preamble for the combined build;
%   * there is NO third copy of the text. The previous
%     ozGUI_ozLib_complete.tex inlined both documents, so every edit had to
%     be made twice and the copies drifted. That file is superseded.
%
% Same pattern as docs/manuscript/manuscript.tex.
%
% Build:  sh build_report.sh          (pdflatex three times: \tableofcontents
%                                      and \part numbering need the extra pass)
% =====================================================================
\documentclass[11pt,a4paper]{article}
%lmodern BEFORE fontenc: Latin Modern ships TS1 (Text Companion) in Type 1.
%Without it, any TS1 glyph -- a degree sign, a micro sign, the upright quotes
%that listings uses -- sends pdflatex looking for a bitmap font it cannot
%generate, and the build dies with
%    !pdfTeX error: pdflatex.exe (file tcrm1000): Font tcrm1000 at 600 not found
%producing NO PDF at all. A full TeX Live installation has the Type 1
%versions and never hits this; a default MiKTeX does.
\usepackage{lmodern}
\usepackage[T1]{fontenc}
\usepackage[a4paper,margin=2.5cm]{geometry}
\usepackage{amsmath,amssymb,amsthm}
\usepackage{graphicx}
%Figures live in docs/figures/. This declaration MUST be here in the master:
%docmute discards each fragment's preamble, so the \graphicspath the two
%fragments carry for their standalone builds never reaches this one. Without
%it every \includegraphics fails with "File not found ... using draft
%setting", which does NOT stop the build -- the PDF is produced with empty
%boxes where the figures should be, so a run can look entirely successful
%and be missing every figure.
%
%The trailing {./} keeps anything still loose in docs/ working, so figures
%can be moved into figures/ without touching a single \includegraphics.
\graphicspath{{figures/}{./}}
\usepackage{booktabs}
\usepackage{multirow}
\usepackage{longtable}
\usepackage{array}
\usepackage{hyperref}
\usepackage{xcolor}
\usepackage{listings}
\usepackage{enumitem}
\usepackage{parskip}
\usepackage{fancyhdr}
\usepackage{titlesec}
\usepackage{comment}

%docmute MUST be loaded after the packages the fragments would otherwise
%load themselves: it discards their preambles, so anything they relied on
%has to be present here. The two fragments currently share this preamble
%exactly, so the union is this list.
\usepackage{docmute}

%LONG IDENTIFIERS RUNNING PAST THE RIGHT MARGIN.
%The report names a lot of code: boltzmannOfP2Ppotential,
%CONSISTENT_PARAMETER_CLOSURES, sasfit_sq_sticky_hard_sphere.c,
%setStarPolymerHighFPotential(). TeX will not hyphenate typewriter text by
%default, so any of these landing near a line end sticks out into the margin.
%A build produced 44 overfull hboxes, the worst 117pt (about 4 cm).
%
%[htt] lets TeX hyphenate inside \texttt, which fixes the identifiers;
%emergencystretch gives it room to stretch a line rather than overrun when
%hyphenation is still not enough. Both are needed: the first alone leaves the
%few cases with no legal break point, the second alone cannot break a long
%word at all.
%
%hyphenat is a standard package and MiKTeX installs it on demand. If you ever
%build on the MSYS2 TeX Live instead, install
%mingw-w64-x86_64-texlive-latex-extra first -- that installation also lacks
%lmodern, which is why the build_report*.log files in this directory record a
%fatal error while MiKTeX builds the same source cleanly.
\usepackage[htt]{hyphenat}
\emergencystretch=2em

%natbib for author-year citations against the SHARED bibliography
%docs/ozgui.bib, which the manuscript also uses. The two fragments used to
%carry a \begin{thebibliography} block each, which meant the report printed
%TWO reference lists and produced duplicate-key warnings wherever both cited
%the same paper (rogersyoung1984 and lee1995 did). One \bibliography here
%replaces both.
%
%The fragments' citation KEYS are unchanged, so no \cite in the text needed
%editing -- tools/bibitem_to_bibtex.py transcribed their \bibitem entries
%into BibTeX keeping the keys. Note ozgui.bib therefore holds nine papers
%twice, under the report's lowercase keys and the manuscript's CamelCase
%ones. That is deliberate: no single document cites both, so each prints a
%correct list, and deduplicating would mean editing 44 \cite commands for no
%gain.
\usepackage[numbers,sort&compress]{natbib}

%The fragments' old thebibliography blocks are commented out with
%\begin{comment} ... \end{comment} rather than deleted, which keeps the
%metadata visible next to the text it belongs to. That needs the `comment`
%package HERE: docmute discards each fragment's preamble, so a
%\usepackage{comment} inside a fragment does not reach the combined build and
%\begin{comment} would fail with "Environment comment undefined".
%
%Note the fragments are then no longer independently buildable unless they
%load `comment' themselves as well -- worth adding to their own preambles too
%if that matters.
\usepackage{comment}

\hypersetup{
  colorlinks=true,
  linkcolor=blue!50!black,
  citecolor=blue!50!black,
  urlcolor=blue!50!black,
  pdftitle={ozGUI / ozLib: Ornstein--Zernike Solver Toolkit},
  pdfauthor={ozGUI/ozLib project}
}

\lstset{
  basicstyle=\ttfamily\small,
  breaklines=true,
  frame=single,
  columns=fullflexible,
  backgroundcolor=\color{gray!8},
  keywordstyle=\color{blue!60!black},
  commentstyle=\color{green!40!black},
  stringstyle=\color{red!60!black},
  language=Python,
  showstringspaces=false
}

\pagestyle{fancy}
\fancyhf{}
\fancyhead[L]{ozGUI / ozLib Documentation}
\fancyhead[R]{\thepage}
\renewcommand{\headrulewidth}{0.4pt}

\newtheorem{remark}{Remark}

\newcommand{\code}[1]{\texttt{#1}}

\title{\bfseries ozGUI / ozLib\\[2mm]
Ornstein--Zernike Solver Toolkit\\[1mm]
\large including the multicomponent and polydisperse extension}
\author{ozGUI/ozLib project}
\date{\today}

\begin{document}
\maketitle
\tableofcontents
\clearpage

%The two fragments each call \maketitle and \tableofcontents inside their own
%document body, because each is also a standalone report. docmute keeps the
%body verbatim, so those calls arrive here too and would print the title and
%contents three times over. Neutralising them HERE rather than editing the
%fragments is what keeps the fragments independently buildable -- which is the
%whole point of this arrangement.
%
%\def, not \renewcommand: the article class does \let\maketitle\relax once
%the master's own \maketitle has run, so \maketitle is then no longer a
%defined command and \renewcommand fails with "Command \maketitle undefined".
\def\maketitle{}
\def\tableofcontents{}

\part{The one-component toolkit}
\input{ozGUI_ozLib_documentation}

\clearpage
\part{Multicomponent and polydisperse extension}
\input{ozGUI_ozLib_multicomponent_supplement}

%ONE bibliography for both parts, replacing the thebibliography block each
%fragment used to carry. Needs a bibtex pass: build_report.sh runs pdflatex,
%bibtex, pdflatex, pdflatex and reports undefined citations, because an
%unresolved \cite renders as (?) in an otherwise clean build.
\clearpage
\bibliography{ozgui}
\bibliographystyle{unsrtnat}

\end{document}
